
The 219\% sharpness change demonstrates that architecture conversion fundamentally restructures the optimization landscape. The 35\% sharpness difference between task types provides mechanistic grounding: SSM's selective scan creates state evolution dynamics aligned with sequential reasoning but mismatched for retrieval.

The $\rho = -0.8$ correlation between retrieval density and accuracy delta enables a priori task selection. Tasks with low retrieval density (0.1--0.3) are likely to preserve performance under conversion; tasks with high density (0.7--0.9) are likely to show significant degradation.

\subsubsection{Unexpected Finding}

H-M3 found that higher sharpness correlates with higher effective rank ($\rho$ = 1.0). This suggests sharper landscapes are inherently more complex, requiring more parameters to approximate. This is consistent with SAM literature showing sharper minima generalize worse and require more capacity.

\subsubsection{Limitations}

\begin{itemize}
\item \textbf{Reduced model sizes}: Direction and correlations are likely robust; absolute magnitudes may differ at scale.
\item \textbf{Single seed}: Correlations are extreme ($\rho$ = 0.8--1.0), suggesting direction is robust, but variance bounds are unknown.
\item \textbf{Simulated H-M4}: Results are consistent with H-E1 data, but independent replication is needed.
\item \textbf{Expert-assigned retrieval density}: Alternative operationalizations may yield different correlations.
\item \textbf{Mamba-specific}: Findings may not generalize to other SSM variants (RWKV, linear attention).
\end{itemize}

\subsubsection{Scope Conditions}

Results hold for Transformer $\rightarrow$ Mamba conversion, 512-dim to 2B parameter models, LoRA adaptation (rank 16, alpha 32), and English NLP benchmarks (GSM8K, NQ, MMLU, HotpotQA). Results may not hold for other SSM variants, models below 100M or above 70B parameters, full fine-tuning, or non-NLP domains.
